% Deutsche Parallelfassung zu why/part_intro.tex

\chapter*{Warum Alignments wesentlich sind}
\phantomsection
\addcontentsline{toc}{chapter}{Einführung}

Die Bearbeitung eines Alignments ist niemals nur die Bearbeitung einer Kurve.
Sie ändert eine konstruktive Definition, erzeugt geometrische Folgen, kann mit
Beobachtungen verglichen werden und kann eine verantwortete
Ingenieurentscheidung erfordern. Dieser Teil untersucht diese Kette, bevor die
später in der Thesis verwendete formale Maschinerie eingeführt wird.

Drei Fragen ordnen die Kapitel: Warum fragmentiert die gegenwärtige Praxis die
Bedeutung eines Alignments? Welchen Mangel muss eine neue Darstellung beheben?
Und welche Art von Framework könnte Ingenieurwissen anwendbar machen, ohne das
Objekt mit einer seiner Beschreibungen zu verwechseln? Die Antworten begründen
die Notwendigkeit von AIM und bereiten die anschließenden mathematischen
Grundlagen vor.
